% =====================================================================
\section{Evidence: humans and animals}
\subsection{What is seen in patients}
\begin{itemize}
\item In the frontal cortex of deceased patients the amounts of parvalbumin and GAD67 mRNA are reduced, while the density of the PV cells themselves is unchanged: the cells are in place but produce less of these proteins.\tP
\item Glutathione. The first study (2000): 27\,\% less in the cerebrospinal fluid, 52\,\% less in the frontal cortex.\tS\ A 2021 synthesis (589 patients and 601 healthy subjects): the mean difference is not significant, but the spread in patients is larger.\tS
\item The capacity to synthesise dopamine is elevated (section~13).\tP
\end{itemize}
\subsection{What has been shown in experiments}
\begin{itemize}
\item \textbf{Ketamine in cell culture.}\tS\ GAD67 and parvalbumin fall by half in 4 and 8 hours; reversibly. This gives the time scale of the slow link and an argument for ``silence'' rather than death.
\item \textbf{Ketamine in mice.}\tS\ One dose, nothing; two consecutive days, a persistent loss. A threshold in duration.
\item \textbf{Receptor blockade in newborn rodents.}\tS\ Glutathione synthesis falls by half in a day. The link ``receptor $\to$ defence''.
\item \textbf{Prenatal brain injury in rats} (the substance MAM).\tS\ Loss of parvalbumin in the hippocampus, then excess dopamine activity in the adults. The bridge to dopamine and the role of adolescence.
\item \textbf{Mice lacking part of the NMDA receptor.}\tS\ A second hit is needed for impairments: a weak receptor shifts the system towards the boundary but does not take it across.
\end{itemize}
\begin{note}
The absence of a mean glutathione deficit in patients does not refute the hypothesis: spectroscopy averages over all cells, and a larger spread may mean a subgroup with a deep deficit or a disturbed daily rhythm. Nor does it confirm it: the explanation ``it is a matter of a subgroup or of a rare cell type'' requires independent verification, otherwise it renders the hypothesis irrefutable.
\end{note}
